% Part: sets-functions-relations
% Chapter: functions
% Section: partial-functions

\documentclass[../../../include/open-logic-section]{subfiles}


\begin{document}

\olfileid[fa-IR]{sfr}{fun}{par}

\olsection{توابع جزئی}

\begin{explain}
گاه سودمند است تعریف تابع را چنان تضعیف کنیم که لازم نباشد
خروجی تابع برای همهٔ ورودی‌های ممکن تعریف شده باشد. چنین
نگاشت‌هایی \emph{توابع جزئی} نامیده می‌شوند.
\end{explain}

\begin{defn}
یک \emph{تابع جزئی} $f \colon A \pto B$ نگاشتی است که
به هر !!{element} از~$A$ حداکثر یک !!{element} از~$B$ نسبت می‌دهد.
اگر $f$ عضوی از~$B$ را به $x \in A$ نسبت دهد، می‌گوییم $f(x)$
\emph{تعریف‌شده} است، و در غیر این صورت \emph{تعریف‌نشده}. اگر $f(x)$ تعریف‌شده باشد،
می‌نویسیم $f(x) \fdefined$، وگرنه $f(x) \fundefined$. \emph{دامنهٔ}
تابع جزئی~$f$ زیرمجموعهٔ~$A$ متشکل از همهٔ عناصری است که تابع روی آن‌ها
تعریف شده است؛ یعنی $\dom{f} = \Setabs{x \in A}{f(x) \fdefined}$.
\end{defn}

\begin{ex}
هر تابع $f\colon A \to B$ یک تابع جزئی نیز هست. توابع
جزئی‌ای که در همه‌جای~$A$ تعریف شده‌اند---یعنی همان توابعی که تاکنون
صرفاً تابع نامیده‌ایم---\emph{تام} نیز نامیده
می‌شوند.
\end{ex}

\begin{ex}
تابع جزئی $f \colon \Real \pto \Real$ که با $f(x) = 1/x$
داده می‌شود، در $x = 0$ تعریف‌نشده و در همه‌جای دیگر تعریف‌شده است.
\end{ex}

\begin{prob}
با داشتن $f\colon A \pto B$، تابع جزئی $g\colon B \pto
A$ را چنین تعریف کنید: برای هر $y \in B$، اگر دقیقاً یک $x \in A$ وجود داشته باشد
چنان‌که $f(x) = y$، آنگاه $g(y) = x$؛ وگرنه $g(y) \fundefined$. نشان دهید
اگر $f$ یک‌به‌یک باشد، آنگاه $g(f(x)) = x$ برای همهٔ $x \in \dom{f}$، و
$f(g(y)) = y$ برای همهٔ $y \in \ran{f}$.
\end{prob}

\begin{defn}[گراف یک تابع جزئی]
فرض کنید $f\colon A \pto B$ تابعی جزئی باشد. \emph{گرافِ}~$f$
رابطهٔ $R_f \subseteq A \times B$ است که به‌صورت زیر تعریف می‌شود:
\[
R_f = \Setabs{\tuple{x,y}}{f(x) = y}.
\]
\end{defn}

\begin{prop}
فرض کنید $R \subseteq A \times B$ این ویژگی را داشته باشد که هرگاه $Rxy$
و $Rxy'$، آنگاه $y = y'$. پس $R$ گراف تابع جزئی
$f\colon A \pto B$ است که چنین تعریف می‌شود: اگر یک $y$ وجود داشته باشد چنان‌که
$Rxy$، آنگاه $f(x) = y$، وگرنه $f(x) \fundefined$. اگر $R$ افزون بر این
\emph{سریال} باشد، یعنی برای هر $x \in A$ یک $y \in B$ وجود داشته باشد چنان‌که
$Rxy$، آنگاه $f$ تام است.
\end{prop}

\begin{proof}
فرض کنید یک $y$ وجود دارد چنان‌که $Rxy$. اگر $y'
\neq y$ دیگری وجود داشت چنان‌که $Rxy'$، شرط مربوط به $R$
نقض می‌شد. بنابراین، اگر یک $y$ وجود داشته باشد چنان‌که $Rxy$، آن $y$
یکتا است و ازاین‌رو $f$ خوش‌تعریف است. آشکارا $R_f = R$ و $f$ در صورتی
تام است که~$R$ سریال باشد.
\end{proof}

\end{document}
